\subsection{Demostración de una cota \texorpdfstring{$\Theta$}{Theta}}
\label{sec:act2-1}

\begin{enunciado}
Demuestre por definición que
\[
  f(n) = 4n^2 \log_2 n + 3n^2 + 7n + 12
\]
pertenece a $\T{n^2 \log_2 n}$. Su demostración debe entregar valores explícitos
$c_1 > 0$, $c_2 > 0$ y $n_0$ tales que, para todo $n \ge n_0$,
\[
  c_1 n^2 \log_2 n \le f(n) \le c_2 n^2 \log_2 n.
\]
\end{enunciado}

Según la definición vista en clases, $f(n) \in \T{n^2 \log_2 n}$ si existen dos
constantes $c_1, c_2 > 0$ y un $n_0 \in \N$ tales que
$0 \le c_1 n^2 \log_2 n \le f(n) \le c_2 n^2 \log_2 n$ para todo $n \ge n_0$.

\textbf{Cota inferior.} Buscamos una constante $c_1$ tal que $c_1 n^2 \log_2 n$ sea
siempre menor o igual que $f(n)$. Para $n \ge 1$, los términos $3n^2$, $7n$ y $12$ son
positivos. Si los quitamos, lo que queda es menor que $f(n)$:
\[
  f(n) = 4n^2 \log_2 n + 3n^2 + 7n + 12 \ge 4n^2 \log_2 n.
\]
Además, $\log_2 n \ge 0$ para $n \ge 1$, así que $4n^2 \log_2 n \ge 0$. Por lo tanto,
$c_1 = 4$ sirve para todo $n \ge 1$.

\textbf{Cota superior.} Buscamos una constante $c_2$ tal que $c_2 n^2 \log_2 n$ sea siempre
mayor o igual que $f(n)$. Para acotar cada término de $f(n)$ por un múltiplo de
$n^2 \log_2 n$ necesitamos que $\log_2 n \ge 1$, y eso ocurre cuando $n \ge 2$.
Con $n \ge 2$:
\begin{itemize}
  \item $4n^2 \log_2 n$ ya es un múltiplo de $n^2 \log_2 n$.
  \item $3n^2 \le 3n^2 \log_2 n$, porque $\log_2 n \ge 1$.
  \item $7n \le 7n^2 \le 7n^2 \log_2 n$, porque $n \le n^2$ y $\log_2 n \ge 1$.
  \item $12 \le 12n^2 \le 12n^2 \log_2 n$, porque $1 \le n^2$ y $\log_2 n \ge 1$.
\end{itemize}
Sumando las cuatro desigualdades:
\[
  f(n) \le 4n^2 \log_2 n + 3n^2 \log_2 n + 7n^2 \log_2 n + 12n^2 \log_2 n
       = 26\, n^2 \log_2 n.
\]
Por lo tanto, $c_2 = 26$ sirve para todo $n \ge 2$.

Juntando ambas cotas, para todo $n \ge 2$:
\[
  4 \cdot n^2 \log_2 n \le f(n) \le 26 \cdot n^2 \log_2 n.
\]
Luego, $c_1 = 4$ y $c_2 = 26$, y las dos cotas valen desde $n = 2$, así que $n_0 = 2$. Por
lo tanto, $f(n) \in \T{n^2 \log_2 n}$.

\noindent\begin{minipage}{\linewidth}
\paragraph{Comprobación.}
Con algunos valores de $n$, las tres cantidades quedan en el orden esperado:
\begin{center}
\begin{tabular}{rrrr}
  \toprule
  $n$ & $4\, n^2 \log_2 n$ & $f(n)$ & $26\, n^2 \log_2 n$ \\
  \midrule
  2 & 16  & 54   & 104  \\
  4 & 128 & 216  & 832  \\
  8 & 768 & 1028 & 4992 \\
  \bottomrule
\end{tabular}
\end{center}
\end{minipage}

